\section{绪论}
\subsection{应用问题与项目边界}
工程应用研究需要从实际问题出发，说明使用者、使用环境、资源约束和预期交付物。研究问题应可以通过设计、实施和测试回答，不能只用器件或软件名称替代问题定义。本工作稿没有实际项目成果，不预设测试结论。
\subsection{已有方案与资料依据}
综述已有工作时，比较需求覆盖、成本、可靠性及适用条件，引用正式出版物、原始技术文档和标准。这里引用可核验著作演示参考文献\cite{zhou2016,lamport1994}，不代表已经完成针对某项工程的技术调研。
\subsection{研究内容与成果形式}
论文、设计和创作可以采用不同的专业组织方式。正式章节应由任务书、专业特点和学院备案规范决定。管理办法列明的正文不少于6000字是正式要求，本示例没有通过重复内容凑够字数，不可直接提交。
\clearpage\section{需求分析与总体设计}
\subsection{需求与验收条件}
表~\ref{tab:trace} 展示需求追踪方式。每项指标应注明来源和验证条件，不能在看到结果后降低验收门槛。涉及安全、隐私或法规的条件应单列，不应与一般性能指标简单平均。
\begin{table}[htbp]\centering\caption{需求追踪表结构示例}\label{tab:trace}
\begin{tabularx}{\linewidth}{lXX}\toprule 需求 & 设计响应 & 验证证据\\\midrule
功能需求 & 定义输入、输出及异常处理 & 功能测试用例与原始记录\\
性能需求 & 明确时间、精度或资源边界 & 在规定条件下重复测量\\
安全需求 & 识别风险并设置限制 & 风险评估与故障测试\\
维护需求 & 提供配置与诊断机制 & 部署及恢复操作记录\\\bottomrule\end{tabularx}
\end{table}
\subsection{系统结构}
图~\ref{fig:architecture} 是一个与具体技术无关的模块图。实际文档应标注接口、通信约束、数据结构和故障路径，而不是只画几个方框。
\begin{figure}[htbp]\centering
\begin{tikzpicture}[node distance=6mm,b/.style={draw,minimum width=65mm,minimum height=10mm,font=\small},>=Stealth]
\node[b](a){输入与边界检查};\node[b,below=of a](b){核心处理与状态管理};
\node[b,below=of b](c){输出、记录及异常反馈};
\draw[->](a)--(b);\draw[->](b)--(c);
\end{tikzpicture}\caption{系统组织结构示例}\label{fig:architecture}
\end{figure}
\subsection{方案比较}
说明候选方案为什么被保留或排除，比较时使用一致的约束和证据。对于无法量化的因素，应明确判断规则，而不是赋予缺乏依据的精确分值。
\clearpage\section{详细实现与质量控制}
\subsection{关键模块与接口}
描述模块的职责、输入、输出和错误处理，并解释设计与需求的对应关系。正文应突出关键实现和决策依据；完整代码、图纸和作品过程材料可放入附录或归档文件，不能把大段源码当作分析。
\subsubsection{参数与配置管理}
把影响结果的参数显式记录，区分开发环境和交付环境。依赖版本、部署方式和运行命令应可复现，涉及授权材料时应记录允许用途和披露范围。
\subsection{指标定义}
若项目需要评价处理时间，可报告平均值以及离散程度。平均值的定义示例如下：
\begin{equation}\bar{t}=\frac{1}{n}\sum_{i=1}^{n}t_i.\end{equation}
式中，$n$ 是满足统计条件的运行次数，$t_i$ 是单次时间。选择指标需考虑异常值、预热状态、硬件环境和并发条件，不能用平均值掩盖严重失败。
\clearpage\section{测试方案与成果评价}
\subsection{测试条件与记录}
测试前确定功能用例、边界条件、故障场景和验收标准。保存原始日志和失败记录，不能只展示成功截图。表~\ref{tab:test} 是记录表框架，不表示任何测试已经通过。
\begin{table}[htbp]\centering\caption{测试记录表框架（未实施）}\label{tab:test}
\begin{tabularx}{\linewidth}{lXXX}\toprule 用例 & 条件与输入 & 预期行为 & 实际记录\\\midrule
正常输入 & 按需求填写 & 按设计填写 & 待实际测试\\
边界输入 & 按需求填写 & 按设计填写 & 待实际测试\\
故障恢复 & 按需求填写 & 按设计填写 & 待实际测试\\\bottomrule\end{tabularx}
\end{table}
\subsection{分析与局限}
将实际结果与预先设定的验收条件比较，说明未达标项目和影响。作品评价或企业反馈应有可核验来源，不得虚构企业采纳证明、使用效果或专家意见。
\section{总结与展望}
总结实际完成的工作，说明研究目标、实现成果和证据之间的对应关系。把未来计划与已完成成果区分开来。本文件只提供论文结构与排版示例，不支持关于工程性能的任何判断。
